\section{Simulation Method}

\subsection{Spin Lattice Initialisation}

The system is modelled on a two-dimensional $(N \times N)$ square lattice with periodic boundary conditions, where each site $(i,j)$ carries a classical spin vector $\mathbf{S}_{i,j}$ constrained to the unit sphere. The initial spin configuration is generated by drawing random orientations uniformly distributed over the unit sphere.

The energy of the system is governed by the spin Hamiltonian containing three contributions: a ferromagnetic Heisenberg exchange interaction with coupling constant $J$, a Dzyaloshinskii--Moriya interaction (DMI) with strength $D$, and a Zeeman coupling to an out-of-plane magnetic field $H_z$. The DMI encodes the form that stabilises Néel-type skyrmions.

\subsection{Monte Carlo Procedure}

The equilibrium spin configuration is sampled using the Metropolis--Hastings algorithm. At each Monte Carlo step, a lattice site is chosen at random and a trial spin is proposed by adding a small random perturbation to each Cartesian component of the current spin, after which the spin is renormalised to the unit sphere. Rather than recomputing the full lattice energy, only the local energy change $\Delta E$ associated with the selected site is evaluated. The trial move is then accepted or rejected according to the Metropolis criterion.

To anneal the system toward its ground state, the simulation is run sequentially at a schedule of increasing inverse temperatures. At each temperature stage the system is thermalised before advancing to the next stage.

\subsection{Energy Stabilisation Check}

Because the simulation starts from a random spin configuration, it must thermalise before meaningful observables can be sampled. Convergence is monitored through a sliding-window criterion on the mean energy. Monte Carlo steps are grouped into windows of size $w = 3N^2$, and at the end of each window the average energy is compared with that of the previous window. The system is considered stable when the relative change between consecutive window averages falls below 15\%. This condition must be satisfied for at least two consecutive windows. Once stability is confirmed, the simulation either advances to the next temperature stage or, at the final stage, begins collecting samples.

\subsection{Skyrmion Detection}

For each sampled spin configuration, the topological charge density $q_{i,j}$ is computed and subsequently convolved with a Gaussian kernel of width $\sigma = 0.8$ lattice spacings, with periodic boundary conditions imposed. This suppresses lattice-scale noise while preserving the broad peaks associated with individual skyrmions.

Skyrmion centres are then identified as local extrema of the smoothed charge density within a neighbourhood of radius 2. An adaptive threshold is applied to suppress spurious detections arising from thermal fluctuations: the threshold fraction is tuned iteratively until the number of detected positive peaks minus negative peaks is consistent with the integer value of $Q$. In this way, weak extrema originating from thermal noise are excluded while the physically relevant skyrmionic structures are retained.
